
\chapter{Parent Theory, Geometry Lift, and Linearized PDE Backbone}
\label{chap:part01_parent_geometry}

\section*{Stage range: 001--023}
\addcontentsline{toc}{section}{Stage range: 001--023}

This chapter is the archive entry that later parts treat as shared infrastructure.  It imports the parent GNLS/Maxwell equations, promotes the throat geometry from the old \((a,L)\) closure to a distributed shape field, and packages the first reduced wall/support/outgoing operators into the grouped notation reused throughout the ledger.

\section*{Part I at a glance}

\begin{readermap}
\item[Primary question] How does the derivation program pass from the parent PDE model to a reusable reduced wall/support/outgoing packet?
\item[Starting inputs] The declared \((4+1)\)-dimensional GNLS/Maxwell parent theory, the old breathing variables \((a,L)\), and the confinement lift \(V_{\rm conf}(\mathbf X;a,L)\to V_{\rm conf}(\mathbf X;\Sigma)\).
\item[Delivers] The moving-surface lift; the breathing reduction; the BdG support Schur complement; the projection-first Maxwell law with its matched one-port normal form; and the grouped \(D,N,P\) normalization language.
\item[Used by] Part II for overlap and spectral placement, Part III for support completion, and Parts IV--VIII for retarded, quotient, realization, and strict/relaxed branch tests.
\item[Result anchors] \resultanchor{MTDC-T1}, \resultanchor{MTDC-T2}, \resultanchor{MTDC-T3}, and \resultanchor{MTDC-T4}; the closing section hands off directly to \resultanchor{MTDC-T5}.
\end{readermap}

\paragraph{Stage packets.}
\begin{center}
\small
\begin{tabular}{@{}>{\raggedright\arraybackslash}p{0.11\textwidth}>{\raggedright\arraybackslash}p{0.24\textwidth}>{\raggedright\arraybackslash}p{0.19\textwidth}>{\raggedright\arraybackslash}p{0.38\textwidth}@{}}
\toprule
Stages & Packet & Status & Carry-forward output \\
\midrule
001--002 & Parent import and geometry lift & \StatusExact{} / \StatusExactClosure{} & Shape field \(R(\Omega,w,t)\), promoted confinement coupling, and the two-mode breathing reduction recovering \((\delta a,\delta L)\) as collective moments. \\
\addlinespace
003 & Conservative BdG support response & \StatusReduced{} / \StatusExactClosure{} & Stable-mode Schur complements, pole shifts, and the first grouped \(P_2\) conservative support diagnostics. \\
\addlinespace
004--021 & Projected EM and parent-action packet & \StatusExact{} / \StatusReduced{} / \StatusExactClosure{} & Projection-first Maxwell leakage, matched mixed-sector normal form, parent-action packet, and the first passive/outgoing \(l=2\) bridge. \\
\addlinespace
022--023 & Grouped normalization packet & \StatusExactClosure{} / \StatusOpen{} & Grouped projectors, normalized \(D,N,P\) coefficients, isotropic ratio \(P_0=N_0/D_0\), and anisotropy transport laws reused later. \\
\bottomrule
\end{tabular}
\end{center}

\claimstatus{Mixed.  The imported parent equations and projection-first Maxwell identities are \StatusExact{} inside the declared parent theory and projection kernel.  The moving-surface algebra, breathing reduction, Schur complements, grouped projectors, and normalized branch formulas are \StatusExactClosure{} once the declared reduced closure is adopted.  Stable-mode reduction, matched Maxwell reduction, low-frequency expansion, and outgoing-port attachment are \StatusReduced{}.  Actual nonlinear branch realization of the final normalization target is \StatusOpen{}.}

\derivationfield{What this part proves}{Part I fixes the parent-theory import, the moving-surface lift, the first reduced wall/BdG response operators, the projection-first Maxwell-to-bundle bridge, and the grouped normalization packet that later parts treat as shared infrastructure.}
\derivationfield{What this part does not prove}{It does not prove that the completed nonlinear moving-throat branch realizes the final outgoing-normalization target; it only derives the reduced object that the branch must realize.}
\derivationfield{Why later parts need it}{Every later part reuses this backbone. Part II builds explicit overlap and selected-mode geometry on top of it, and all later closure, quotient, and realization arguments inherit its operator and normalization language.}

\section{Block contract and theorem-level output}
\label{sec:part01-block-contract}

This part has a narrow verification contract.  It must show that the later reduced response program is not an unrelated collection of coefficients.  The coefficients used later arise from one derivation chain:
\begin{equation}
\label{eq:part01-chain-overview}
\begin{aligned}
\text{parent GNLS/Maxwell theory}
&\longrightarrow
\text{moving shape field}
\longrightarrow
\text{linear wall/BdG/Maxwell system}
\\
&\longrightarrow
\text{grouped response bundle}.
\end{aligned}
\end{equation}
The chain is deliberately status-stratified.  The parent action and its Euler--Lagrange identities are exact inside the declared parent model.  The shape-field lift is the working effective geometry closure.  The projected Maxwell law is exact once the observation kernel is declared, while the old reduction-first Maxwell block survives only as a matched one-port normal form.  The stable-mode and Maxwell/mixed Schur complements are exact after their reduced normal-coordinate ansatz is declared.  The final outgoing normalization condition is an exact branch test, but its realization by the actual moving-throat branch remains open at this stage.

\begin{theorem}[Part I backbone theorem]
\label{thm:part01-backbone}
Under the declared parent \((4+1)\)-dimensional gauged GNLS plus localized Maxwell theory, and under the moving-surface closure
\begin{equation}
\label{eq:part01-shape-field-short}
\Sigma(\mathbf X,t)=r-R(\Omega,w,t),
\qquad
r=\sqrt{x^2+y^2+z^2},
\qquad
\Omega=\mathbf x/r,
\end{equation}
the linearized throat problem admits a reduced wall/support/outgoing response description with the following structure.
\begin{enumerate}[leftmargin=2.2em]
  \item The old variables \(a(t)\) and \(L(t)\) are recovered as lowest collective moments of \(R\), not treated as fundamental moving-throat coordinates.
  \item The real \(l=2\) wall/support sector splits into grouped lanes \(20,21,22\), with the \(21\) and \(22\) entries denoting real cosine/sine pairs.
  \item Stable matter support modes generate conservative Schur-complement corrections to the wall operator.
  \item Projection-first localized Maxwell theory gives an open-system brane law; its matched mixed \(A_w/F_{\mu w}/J^w\) normal form generates conservative self-energies and a passive/outgoing transfer factor.
  \item On the isotropic grouped branch, the leading outgoing quadrupole test reduces to the single invariant ratio
  \begin{equation}
  \label{eq:part01-normalization-product-overview}
  \widehat m_0^{\,2}P_0
  =
  \frac{54G c_s^5}{5a^5c^5},
  \qquad
  P_0=\frac{N_0}{D_0}.
  \end{equation}
\end{enumerate}
The derivation of \eqref{eq:part01-normalization-product-overview} is exact inside the reduced grouped response closure.  Whether the actual nonlinear moving-throat branch realizes the required value is not settled in this part.
\end{theorem}

\begin{proof}
The statement is a compilation of the twenty-three stage-level reductions.  The moment recovery and modal decomposition follow from the shape-field definition.  The breathing reduction follows by inserting the lowest axisymmetric ansatz into the quadratic wall action.  The BdG kernel follows by an exact Schur complement in the declared stable-mode coordinates.  The Maxwell sector first follows by exact projection of the localized parent equation, then by matched mouth/bundle transport, parent-action packet bookkeeping, and the reduced Schur complement for the one-port mixed normal form.  The final normalization product follows by converting the conservative grouped operator into a normalized response and multiplying by the compact outgoing \(l=2\) fingerprint.  Each individual step is written below with its own status boundary.
\end{proof}

\section{Exact parent action and field-equation import}
\label{sec:part01_parent_geometry:exact-parent-action-and-field-equation-import}

\subsection{Local goal and import boundary}
\label{subsec:part01-parent-goal}

The moving-throat program does not begin by inventing new field equations.  It imports the already-declared parent \((4+1)\)-dimensional matter/gauge theory and then promotes the old geometry closure into a distributed moving-surface variable.  The imported parent equations are therefore the fixed background against which the moving-throat extension must be checked.

The exact import contains three pieces:
\begin{enumerate}[leftmargin=2.2em]
\item a gauged four-spatial-dimensional GNLS matter sector;
\item a localized \((4+1)\)-dimensional Maxwell sector with localization profile \(Z(w)\);
\item a confinement coupling that is later promoted from \(V_{\rm conf}(\mathbf X;a,L)\) to \(V_{\rm conf}(\mathbf X;\Sigma)\).
\end{enumerate}
The geometry lift can modify the confinement argument, the wall response, and the branch data.  It does not modify the charge ontology, the minimal-coupling rule, or the exact localized Maxwell equation.

\subsection{Arena and core fields}
\label{subsec:part01-parent-arena}

The parent arena has coordinates
\begin{equation}
\label{eq:part01-parent-coordinates}
x^M=(t,x,y,z,w),
\qquad
M,N\in\{0,1,2,3,4\}.
\end{equation}
Bulk spatial coordinates are
\begin{equation}
\mathbf X=(x,y,z,w)\in\mathbb R^4,
\qquad
\mathbf x=(x,y,z)\in\mathbb R^3,
\end{equation}
where \(\mathbf x\) is the brane-facing spatial coordinate.  The flat bulk metric convention is
\begin{equation}
\label{eq:part01-metric}
\eta_{MN}=\operatorname{diag}(-1,+1,+1,+1,+1),
\qquad
\Box_5=-\partial_t^2+\nabla_3^2+\partial_w^2.
\end{equation}
The core parent fields are
\begin{equation}
\psi(\mathbf X,t),
\qquad
A_M(x),
\qquad
F_{MN}=\partial_MA_N-\partial_NA_M,
\qquad
\rho=|\psi|^2.
\end{equation}
The moving-throat extension introduces either a level set \(\Sigma(\mathbf X,t)\) or the equivalent shape field \(R(\Omega,w,t)\), but that variable enters through the promoted confinement coupling.

\subsection{Charge ontology imported into the PDE program}
\label{subsec:part01-charge-ontology-import}

The electric branch label is
\begin{equation}
\label{eq:part01-charge-ontology}
\eta_Q=\pm1,
\qquad
q_\star=\eta_Q e_\star,
\qquad
e_\star>0.
\end{equation}
After zero-mode brane normalization, the observed brane charge is
\begin{equation}
\label{eq:part01-effective-charge}
q_{\rm eff}=\frac{q_\star}{\sqrt{Z_{\rm int}}},
\qquad
e_{\rm eff}=\frac{e_\star}{\sqrt{Z_{\rm int}}},
\qquad
Z_{\rm int}=\int_{-\infty}^{+\infty}Z(w)\,dw.
\end{equation}
This firewall matters already in Part I because the mixed channels \(A_w,J^w,F_{\mu w}\) are retained microscopically.  Keeping those channels alive must not be confused with returning to an older circulation-based definition of electric charge.  Circulation remains a magnetic/vortical label; it is not the sign or magnitude of the electric branch.

\subsection{Matter action and exact GNLS equation}
\label{subsec:part01-matter-action}

The imported matter Lagrangian density is
\begin{equation}
\label{eq:part01-Lpsi}
\mathcal L_\psi
=
\frac{i\hbar}{2}\left(\psi^\ast D_t\psi-\psi D_t\psi^\ast\right)
-\frac{\hbar^2}{2m}(D_i\psi)^\ast(D_i\psi)
-V_{\rm conf}(\mathbf X;\Sigma)\rho
-U(\rho).
\end{equation}
At the exact parent level before promotion, the confinement potential is usually written as \(V_{\rm conf}(\mathbf X;a,L)\).  In the moving-throat lift the same coupling is read as a surface-dependent potential.  The covariant derivatives are
\begin{equation}
\label{eq:part01-covariant-derivatives}
D_t\psi=\partial_t\psi+\frac{iq_\star}{\hbar}A_0\psi,
\qquad
D_i\psi=\partial_i\psi-\frac{iq_\star}{\hbar}A_i\psi.
\end{equation}
The stiff-polytropic equation of state carried into the program is
\begin{equation}
\label{eq:part01-eos}
P(\rho)=K\rho^5,
\qquad
U(\rho)=\frac K4\rho^5,
\qquad
h(\rho)=\frac{dU}{d\rho}=\frac{5K}{4}\rho^4,
\end{equation}
with sound-speed ledger
\begin{equation}
\label{eq:part01-sound-speed}
c_s^2(\rho)=\frac1m\frac{dP}{d\rho}=\frac{5K}{m}\rho^4.
\end{equation}
Variation with respect to \(\psi^\ast\) gives the exact GNLS equation
\begin{equation}
\label{eq:part01-gnls}
\boxed{
i\hbar D_t\psi
=
\left[-\frac{\hbar^2}{2m}D_iD_i+V_{\rm conf}(\mathbf X;\Sigma)+h(\rho)\right]\psi}.
\end{equation}
The exact number current and continuity equation are
\begin{equation}
\label{eq:part01-number-current}
\boxed{
j^i=\frac{\hbar}{m}\Im(\psi^\ast D_i\psi),
\qquad
\partial_t\rho+\partial_i j^i=0}.
\end{equation}

\subsection{Localized Maxwell action and exact mixed observables}
\label{subsec:part01-maxwell-action}

The localized Maxwell Lagrangian density is
\begin{equation}
\label{eq:part01-LEM}
\mathcal L_{\rm EM}
=
-\frac{Z(w)}{4\mu_0}F_{MN}F^{MN}
-\frac{H(w)}{2\xi\mu_0}(\partial\!\cdot\!A)^2
-A_MJ_{\rm ext}^M.
\end{equation}
Here \(H(w)=1\) is the old unweighted gauge-fixing convention, while
\(H(w)=Z(w)\) is the projection-stable localized convention used below.
The matter current generated by minimal coupling is not to be double-counted in \(J_{\rm ext}^M\).  The source entering Maxwell's equation is
\begin{equation}
J_{\rm tot}^M=J_\psi^M+J_{\rm ext}^M.
\end{equation}
Variation with respect to \(A_N\) gives
\begin{equation}
\label{eq:part01-maxwell-equation}
\boxed{
\partial_M\!\bigl(Z(w)F^{MN}\bigr)
+\frac1\xi\partial^N\!\bigl(H(w)\partial\!\cdot\!A\bigr)
=
\mu_0J_{\rm tot}^N}.
\end{equation}
The Bianchi identity is
\begin{equation}
\partial_{[L}F_{MN]}=0.
\end{equation}
The mixed gauge-invariant fields retained by the PDE program are
\begin{equation}
\label{eq:part01-mixed-fields}
\boxed{
E_w=F_{w0}=-\partial_tA_w-\partial_wA_0,
\qquad
C_a=F_{aw}=\partial_aA_w-\partial_wA_a}.
\end{equation}
They are suppressed only in the strict far-field brane Maxwell reduction.  They remain available in the moving-throat PDE and become essential in the first outgoing bridge.

\subsection{Madelung identities and the vortex/gauge firewall}
\label{subsec:part01-madelung-import}

Writing
\begin{equation}
\psi=\sqrt\rho\,e^{i\theta},
\end{equation}
the exact gauge-invariant velocity is
\begin{equation}
\label{eq:part01-madelung-velocity}
v_i=\frac{\hbar}{m}\partial_i\theta-\frac{q_\star}{m}A_i.
\end{equation}
The quantum potential is
\begin{equation}
Q(\rho)=-\frac{\hbar^2}{2m}\frac{\nabla_4^2\sqrt\rho}{\sqrt\rho}.
\end{equation}
The Euler-like form is
\begin{equation}
\label{eq:part01-euler-form}
m(\partial_t+v_j\partial_j)v_i
=
q_\star(E_i+v_jB_{ij})
-\partial_i\left(V_{\rm conf}+h(\rho)+Q(\rho)\right),
\end{equation}
and the exact vorticity--gauge identity is
\begin{equation}
\label{eq:part01-vorticity-gauge}
\boxed{
\Omega_{ij}:=\partial_i v_j-\partial_jv_i=-\frac{q_\star}{m}F_{ij}}.
\end{equation}
This equation is frequently useful, but it is not a charge dictionary.  It records the relation between gauge field strength and gauge-invariant vorticity away from singular phase defects.

\subsection{Projection and reduction hooks}
\label{subsec:part01-projection-hooks}

For a normalized brane-observation kernel \(W(w)\), projection means
\begin{equation}
\label{eq:part01-projection-operator}
\langle Q\rangle_W=\int_D W(w)Q(w)\,dw,
\qquad
\int_DW(w)\,dw=1.
\end{equation}
Projection commutes with brane derivatives, but not with the transverse
derivative:
\begin{equation}
\label{eq:part01-projection-ibp}
\langle \partial_w Q\rangle_W
=\bigl[WQ\bigr]_{\partial D}-\int_DW'(w)Q(w)\,dw.
\end{equation}
This identity is the reason a projection-first brane theory is generally an
open system.

Applied to the matter current, the same projection defines
\begin{equation}
\rho_{\rm brane}(\mathbf x,t)=\int W(w)\rho(\mathbf x,w,t)\,dw,
\qquad
\mathbf j_{\rm brane}=\int W(w)\mathbf j_{xyz}(\mathbf x,w,t)\,dw.
\end{equation}
Projected continuity is exact once \(W\) is declared:
\begin{equation}
\label{eq:part01-projected-continuity}
\partial_t\rho_{\rm brane}+\nabla_3\cdot\mathbf j_{\rm brane}=S_{\rm leak},
\end{equation}
where
\begin{equation}
\label{eq:part01-leakage-source}
S_{\rm leak}
=-[Wj^w]_{-\infty}^{+\infty}+\int W'(w)j^w\,dw.
\end{equation}

Applied to the localized Maxwell equation \eqref{eq:part01-maxwell-equation},
projection gives, for brane index \(\nu\),
\begin{equation}
\label{eq:part01-projected-maxwell-law}
\boxed{
\partial_\mu\langle ZF^{\mu\nu}\rangle_W
+\bigl[WZF^{w\nu}\bigr]_{\partial D}
-\int_DW'(w)ZF^{w\nu}\,dw
+\frac1\xi\partial^\nu\langle H\,\partial\!\cdot\!A\rangle_W
=\mu_0\langle J^\nu\rangle_W}.
\end{equation}
The homogeneous equations project in the measured fields
\begin{equation}
\label{eq:part01-measured-vs-flux-fields}
E_{{\rm meas},a}=\langle F_{a0}\rangle_W,
\qquad
B_{{\rm meas},ab}=\langle F_{ab}\rangle_W,
\end{equation}
whereas \eqref{eq:part01-projected-maxwell-law} uses the source-coupled fluxes
\begin{equation}
D_{\rm flux}^{a}=\langle ZF^{a0}\rangle_W,
\qquad
H_{\rm flux}^{ab}=\langle ZF^{ab}\rangle_W.
\end{equation}
Identifying measured fields with flux fields is therefore a closure assumption,
not a parent identity.

In the zero-mode channel
\begin{equation}
A_\mu(x,w)=a_\mu(x),
\qquad
A_w=0,
\qquad
F^{w\nu}=0,
\qquad
J^\nu(x,w)=j^\nu(x)S(w),
\end{equation}
the projected effective parameters are
\begin{equation}
\label{eq:part01-projected-effective-parameters}
\boxed{
\mu_{0,\rm eff}^{\rm proj}
=\mu_0\frac{\int WS\,dw}{\int WZ\,dw},
\qquad
\xi_{\rm eff}^{\rm proj}
=\xi\frac{\int WZ\,dw}{\int WH\,dw}}.
\end{equation}
The matched reduction-first brane Maxwell law is recovered by the localized
gauge choice and source profile
\begin{equation}
\label{eq:part01-projection-reduction-match}
\boxed{
H=Z,
\qquad
S(w)=\frac{Z(w)}{Z_{\rm int}},
\qquad
\mu_{0,\rm eff}^{\rm proj}=\frac{\mu_0}{Z_{\rm int}}}.
\end{equation}
Thus the reduction-first Maxwell sector is retained below as a controlled
matched channel and one-port normal form, not as the generic parent projection.

Under the Helmholtz split
\begin{equation}
\mathbf v_{\rm brane}=\nabla_3\varphi+\mathbf v_T,
\qquad
\nabla_3\cdot\mathbf v_T=0,
\end{equation}
the exact longitudinal identity is
\begin{equation}
\label{eq:part01-longitudinal-identity}
\rho_{\rm brane}\nabla_3^2\varphi
=
S_{\rm leak}-\partial_t\rho_{\rm brane}
-(\nabla_3\rho_{\rm brane})\cdot(\nabla_3\varphi+\mathbf v_T).
\end{equation}
This is an imported lower-order hook rather than a new Stage~001--023 result.  It fixes how the moving-throat response must later interface with the brane-facing scalar sector.

\begin{theorem}[Parent-theory import]
\label{thm:part01-parent-import}
Within the corrected charge ontology and the declared localized Maxwell/GNLS parent action, the moving-throat derivation inherits the exact equations \eqref{eq:part01-gnls}, \eqref{eq:part01-number-current}, \eqref{eq:part01-maxwell-equation}, \eqref{eq:part01-mixed-fields}, and \eqref{eq:part01-vorticity-gauge}.  After a brane-observation kernel is declared, it also inherits the exact projected Maxwell law \eqref{eq:part01-projected-maxwell-law}.  Any moving-throat closure must therefore preserve the mixed-channel ontology and may only suppress \(A_w,J^w,F_{\mu w},E_w,C_a\) as part of a controlled matched brane reduction.
\end{theorem}

\begin{proof}
The equations are direct Euler--Lagrange consequences of the imported parent action and the minimal-coupling definitions.  The projected Maxwell law follows by applying \eqref{eq:part01-projection-ibp} to the transverse derivative in \eqref{eq:part01-maxwell-equation}.  The mixed fields in \eqref{eq:part01-mixed-fields} are components of the gauge-invariant field strength.  Therefore they cannot be removed by a gauge choice.  The zero-mode brane Maxwell reduction may set them to zero as a matched branch assumption, but the moving-throat PDE that is supposed to compute microscopic response data must retain them until such a branch restriction is explicitly imposed.
\end{proof}

\section{Hybrid level-set / shape-field representation}
\label{sec:part01_parent_geometry:hybrid-level-set-shape-field-representation}

\subsection{Purpose of the geometry lift}
\label{subsec:part01-geometry-lift-purpose}

The old conservative hierarchy used finite-dimensional geometry variables \(a(t)\) and \(L(t)\).  Those variables are adequate as a reduced breathing closure, but they cannot expose the grouped real \(P_2\) support lanes or the full mouth/worldtube response operator.  Stage~001 therefore replaces the old finite-dimensional closure by a distributed surface field while preserving \(a\) and \(L\) as low moments.

The chosen representation is hybrid: the finite throat is represented by a level set, but the level set is parameterized by a brane-spherical mouth shape.  Define
\begin{equation}
\label{eq:part01-hybrid-lift}
\boxed{
\Sigma(\mathbf X,t)=r-R(\Omega,w,t),
\qquad
r=\sqrt{x^2+y^2+z^2},
\qquad
\Omega=\frac{\mathbf x}{r}\in S^2}.
\end{equation}
The throat surface is \(\Sigma(\mathbf X,t)=0\).  We use the sign convention
\begin{equation}
\Sigma>0 \quad\hbox{outside the throat},
\qquad
\Sigma<0 \quad\hbox{inside the support region}.
\end{equation}

\subsection{Reference stationary throat}
\label{subsec:part01-reference-throat}

A stationary finite-throat reference branch is written
\begin{equation}
\label{eq:part01-stationary-surface}
\Sigma_0(\mathbf X)=r-R_0(w).
\end{equation}
The reference mouth is at \(w=0\), the throat extends over \(0\le w\le L_0\), and the reference mouth radius is
\begin{equation}
\label{eq:part01-reference-mouth-radius}
a_0=R_0(0).
\end{equation}
The bottom may be represented by
\begin{equation}
R_0(L_0)=0,
\end{equation}
or by an equivalent regular bottom condition.  The finite-throat assumption is important because the later support problem uses finite axial spectral data rather than a noncompact Gaussian surrogate.

\subsection{Promoted confinement potential}
\label{subsec:part01-promoted-confinement}

The old confinement potential is promoted by replacing
\begin{equation}
V_{\rm conf}(\mathbf X;a,L)
\quad\longrightarrow\quad
V_{\rm conf}(\mathbf X;\Sigma).
\end{equation}
The minimal smooth-wall choice is
\begin{equation}
\label{eq:part01-promoted-wall-potential}
V_{\rm conf}(\mathbf X;\Sigma)
=
V_{\rm wall}\!\left(\frac{\Sigma(\mathbf X,t)}{\ell_c}\right),
\end{equation}
where \(\ell_c\) is a wall-thickness scale.  Around \(\Sigma_0\), the first variation is
\begin{equation}
\delta V_{\rm conf}
=
\frac{V_{\rm wall}'(\Sigma_0/\ell_c)}{\ell_c}\,\delta\Sigma.
\end{equation}
Since \(\Sigma=r-R\), one has
\begin{equation}
\label{eq:part01-delta-sigma}
\delta\Sigma=-\delta R.
\end{equation}
Writing the wall displacement as \(\eta=\delta R\), the moving-wall coupling to the matter sector is therefore
\begin{equation}
\label{eq:part01-delta-V-conf}
\boxed{
\delta V_{\rm conf}
=
-\frac{V_{\rm wall}'(\Sigma_0/\ell_c)}{\ell_c}\,\eta}.
\end{equation}
This sign is one of the main sign checks for the entire linearized program: a positive outward radial displacement decreases \(\Sigma\) at fixed bulk point and therefore enters the confinement variation with the minus sign in \eqref{eq:part01-delta-V-conf}.

\section{Collective moment recovery for \texorpdfstring{\(a(t)\)}{a(t)} and \texorpdfstring{\(L(t)\)}{L(t)}}
\label{sec:part01_parent_geometry:collective-moment-recovery-for-a-t-and-l-t}

\subsection{Mouth-radius and depth moments}
\label{subsec:part01-collective-moments}

The old mouth radius is not discarded.  It is recovered as the spherical average of the moving mouth:
\begin{equation}
\label{eq:part01-a-moment}
\boxed{
a(t)=\frac{1}{4\pi}\int_{S^2}R(\Omega,0,t)\,d\Omega}.
\end{equation}
Let \(W_b(\Omega,t)\) denote the bottom location, implicitly defined by
\begin{equation}
R(\Omega,W_b(\Omega,t),t)=0.
\end{equation}
Then the old axial depth is recovered as
\begin{equation}
\label{eq:part01-L-moment}
\boxed{
L(t)=\frac{1}{4\pi}\int_{S^2}W_b(\Omega,t)\,d\Omega}.
\end{equation}
Equations \eqref{eq:part01-a-moment} and \eqref{eq:part01-L-moment} are the first audit gate for the lift: the new distributed theory must reduce to the old \((a,L)\) closure when restricted to its lowest axisymmetric moments.

\subsection{Mode decomposition of the throat shape}
\label{subsec:part01-mode-decomposition}

Write
\begin{equation}
\label{eq:part01-R-eta}
R(\Omega,w,t)=R_0(w)+\eta(\Omega,w,t).
\end{equation}
The wall displacement is expanded in real spherical harmonics:
\begin{equation}
\label{eq:part01-eta-decomposition}
\eta(\Omega,w,t)
=
q_{00}(w,t)Y_{00}(\Omega)
+
\sum_{A\in P_2^{\rm real}}q_A(w,t)Y_A(\Omega)
+
\eta_{\ge3}(\Omega,w,t).
\end{equation}
Here
\begin{equation}
P_2^{\rm real}=\{20,21c,21s,22c,22s\}.
\end{equation}
The grouped real ledger later compresses the \(c/s\) pairs into grouped labels
\begin{equation}
A\in\{20,21,22\},
\end{equation}
with natural weights \((1,2,2)\).  The grouped labels are lane labels, not spacetime indices.

With the standard normalized harmonic
\begin{equation}
\label{eq:part01-Y00}
Y_{00}=\frac{1}{2\sqrt\pi},
\qquad
\int_{S^2}Y_{00}^2\,d\Omega=1,
\qquad
\frac1{4\pi}\int_{S^2}Y_{00}\,d\Omega=\frac1{2\sqrt\pi},
\end{equation}
the physical mouth-average shift obeys
\begin{equation}
\label{eq:part01-q00-deltaa}
\boxed{
q_{00}(0,t)=2\sqrt\pi\,\delta a(t)}.
\end{equation}
All grouped real \(P_2\) harmonics have zero average over the sphere.  Thus grouped quadrupole deformations do not change the mouth-average radius at linear order.

\subsection{Axisymmetric residual lane}
\label{subsec:part01-axisymmetric-residual-lane}

The axisymmetric wall field is further split as
\begin{equation}
\label{eq:part01-axisymmetric-split}
q_{00}(w,t)Y_{00}(\Omega)
=
2\sqrt\pi\left[\alpha_a(w)\delta a(t)+\alpha_L(w)\delta L(t)\right]Y_{00}(\Omega)
+
g(w,t)Y_{00}(\Omega),
\end{equation}
where \(g(w,t)\) is the residual axisymmetric geometry lane orthogonal to the chosen \(a\) and \(L\) profiles.  In practice, the first truncation drops \(g\); later geometry-completion terms are naturally interpreted as the return of this residual lane.

\begin{proposition}[Collective-moment recovery]
\label{prop:part01-collective-recovery}
Inside the shape-field closure \(\Sigma=r-R(\Omega,w,t)\), the old collective variables \(a(t)\) and \(L(t)\) are recovered as the monopole moments \eqref{eq:part01-a-moment} and \eqref{eq:part01-L-moment}.  The grouped real \(P_2\) deformations are orthogonal to the mouth-average radius shift at linear order.
\end{proposition}

\begin{proof}
The formula for \(a(t)\) is the spherical average of the mouth radius.  The formula for \(L(t)\) is the spherical average of the bottom location.  The zero-average property of all \(l=2\) spherical harmonics gives the last statement immediately.  The normalization bridge \eqref{eq:part01-q00-deltaa} follows from \eqref{eq:part01-Y00}.
\end{proof}

\section{Distributed wall action and densitized convention}
\label{sec:part01_parent_geometry:distributed-wall-action-and-densitized-convention}

\subsection{Quadratic wall action}
\label{subsec:part01-quadratic-wall-action}

The first distributed wall action is the minimal passive local quadratic action
\begin{equation}
\label{eq:part01-S-eta-raw}
S_\eta^{(2)}
=
\frac12\int dt\,dw\,d\Omega\,\sqrt{\gamma_0}\left[
\mu_\eta(w)(\partial_t\eta)^2
-T_w(w)(\partial_w\eta)^2
-T_\Omega(w)\eta(-\Delta_{S^2})\eta
-K_\eta(w)\eta^2
\right].
\end{equation}
The coefficients \(\mu_\eta,T_w,T_\Omega,K_\eta\) are branch data of the chosen reference throat.  They are not to be refit independently at later stages to force a normalization condition.

After integrating over the reference sphere, the program uses a densitized one-dimensional convention: the background surface weight is absorbed into the effective axial coefficients.  With that convention, the modal action is written with measure \(dw\).  If an explicit axial weight \(g(w)=\sqrt{\gamma_0}\) were retained, the Euler--Lagrange equation would contain the corresponding first-derivative contribution from \(g'(w)\).  In the ledger formulas below, that bookkeeping has already been absorbed into \(\mu_\eta,T_w,T_\Omega,K_\eta\).

\subsection{Modal wall equation}
\label{subsec:part01-modal-wall-equation}

For a harmonic component \(q_{lm}(w,t)Y_{lm}(\Omega)\), using
\begin{equation}
-\Delta_{S^2}Y_{lm}=l(l+1)Y_{lm},
\end{equation}
the densitized modal wall equation is
\begin{equation}
\label{eq:part01-modal-wall-equation}
\boxed{
\mu_\eta q_{lm,tt}
-
\partial_w\!\bigl(T_w\partial_w q_{lm}\bigr)
+
\bigl[K_\eta+l(l+1)T_\Omega\bigr]q_{lm}
=
S_{lm}^{(\psi,A)}+f_{lm}^{\rm ext}}.
\end{equation}
The source \(S_{lm}^{(\psi,A)}\) is generated by variation of the matter and gauge sectors, with the first matter contribution inherited from \eqref{eq:part01-delta-V-conf}.

Equation \eqref{eq:part01-modal-wall-equation} separates the first three relevant sectors:
\begin{enumerate}[leftmargin=2.2em]
\item the \(l=0\) collective sector containing \(\delta a\), \(\delta L\), and the residual geometry lane \(g(w,t)\);
\item the grouped real \(l=2\) sector containing \(20,21c,21s,22c,22s\);
\item higher \(l\ge3\) wall modes.
\end{enumerate}
The present response program keeps higher modes in the full PDE ontology but does not need them for the first grouped-quadrupole bridge.

\begin{proposition}[Wall modal backbone]
\label{prop:part01-wall-modal-backbone}
The quadratic wall action \eqref{eq:part01-S-eta-raw}, read in the densitized axial convention, yields the modal operator in \eqref{eq:part01-modal-wall-equation}.  The scalar and grouped real \(P_2\) sectors therefore split before any matter or gauge coupling is added.
\end{proposition}

\begin{proof}
Insert \(\eta=q_{lm}(w,t)Y_{lm}(\Omega)\) into \eqref{eq:part01-S-eta-raw}, use orthonormality on \(S^2\), and vary the resulting one-dimensional action.  The angular stiffness term contributes \(l(l+1)T_\Omega q_{lm}\).  The remaining source terms are not part of the free wall action and are kept on the right-hand side.
\end{proof}

\section{Axisymmetric breathing reduction}
\label{sec:part01_parent_geometry:axisymmetric-breathing-reduction}

\subsection{Two-mode breathing ansatz}
\label{subsec:part01-breathing-ansatz}

Stage~002 checks that the distributed wall lift recovers the old \((a,L)\) closure at the expected conservative linear level.  Use the axisymmetric truncation
\begin{equation}
\label{eq:part01-breathing-ansatz}
\eta_0(w,t)
=
2\sqrt\pi\left[\alpha_a(w)\delta a(t)+\alpha_L(w)\delta L(t)\right].
\end{equation}
Define
\begin{equation}
Q^A=(\delta a,\delta L),
\qquad
A\in\{a,L\}.
\end{equation}
The truncation is not the final geometry theory.  It is the lowest-mode audit that the distributed wall has not lost contact with the older closure.

\subsection{Reduced matrices}
\label{subsec:part01-breathing-matrices}

For the axisymmetric branch, the reduced Lagrangian is
\begin{equation}
\label{eq:part01-breathing-Lred}
L_{\rm red}^{(0)}
=
\frac12M_{AB}\dot Q^A\dot Q^B
-
\frac12K_{AB}Q^AQ^B,
\end{equation}
where
\begin{equation}
\label{eq:part01-breathing-M}
\boxed{
M_{AB}
=
4\pi\int dw\,\mu_\eta(w)\alpha_A(w)\alpha_B(w)},
\end{equation}
and
\begin{equation}
\label{eq:part01-breathing-K}
\boxed{
K_{AB}
=
4\pi\int dw\left[
T_w(w)\alpha_A'(w)\alpha_B'(w)
+
K_0(w)\alpha_A(w)\alpha_B(w)
\right]}.
\end{equation}
The Euler--Lagrange equation is
\begin{equation}
\label{eq:part01-breathing-EOM}
\boxed{
M_{AB}\ddot Q^B+K_{AB}Q^B=0}.
\end{equation}
This is the conservative part of the older effective geometry law.  Damping and open-system terms can be added only after coupling to exterior, matter, gauge, or leakage channels.

\subsection{One-mode grouped \texorpdfstring{\(P_2\)}{P2} comparison}
\label{subsec:part01-p2-one-mode-comparison}

For a single grouped real quadrupole amplitude,
\begin{equation}
\eta_{2m}(\Omega,w,t)=\beta_2(w)q_{2m}(t)Y_{2m}^{\rm real}(\Omega),
\end{equation}
the reduced one-mode Lagrangian is
\begin{equation}
L_{2m}
=
\frac12M_2\dot q_{2m}^{\,2}-\frac12K_2q_{2m}^{\,2},
\end{equation}
with
\begin{equation}
\label{eq:part01-M2-K2}
\boxed{
M_2=\int dw\,\mu_\eta(w)\beta_2(w)^2,
\qquad
K_2=\int dw\left[T_w(w)\beta_2'(w)^2+\bigl(K_\eta(w)+6T_\Omega(w)\bigr)\beta_2(w)^2\right]}.
\end{equation}
Therefore
\begin{equation}
M_2\ddot q_{2m}+K_2q_{2m}=0.
\end{equation}
Before anisotropy or support coupling is added, every real \(P_2\) component has the same geometry-only operator.  This is the first degeneracy statement behind the later grouped-isotropy gates.

\begin{theorem}[Breathing reduction as the lowest wall truncation]
\label{thm:part01-breathing-reduction}
Inside the quadratic distributed-wall closure, the two-mode axisymmetric ansatz \eqref{eq:part01-breathing-ansatz} reduces exactly to the conservative matrix system \eqref{eq:part01-breathing-EOM}.  The old \((a,L)\) variables are therefore the lowest collective coordinates of the wall PDE rather than independent microscopic primitives.
\end{theorem}

\begin{proof}
Substitute \eqref{eq:part01-breathing-ansatz} into the quadratic wall action, use the normalization of \(Y_{00}\), and collect coefficients of \(\dot Q^A\dot Q^B\) and \(Q^AQ^B\).  The resulting matrices are \eqref{eq:part01-breathing-M} and \eqref{eq:part01-breathing-K}.  Varying \eqref{eq:part01-breathing-Lred} gives \eqref{eq:part01-breathing-EOM}.
\end{proof}

\section{Stable BdG normal-mode reduction and Schur complement}
\label{sec:part01_parent_geometry:stable-bdg-normal-mode-reduction-and-schur-complem}

\subsection{Linearized matter coupling}
\label{subsec:part01-linearized-matter-coupling}

Let the stationary matter background be
\begin{equation}
\psi(\mathbf X,t)=e^{-i\mu_0t/\hbar}\psi_0(\mathbf X),
\end{equation}
and write
\begin{equation}
\psi=e^{-i\mu_0t/\hbar}\left(\psi_0+\delta\psi\right).
\end{equation}
The linearized density is
\begin{equation}
\delta\rho=\psi_0^\ast\delta\psi+\psi_0\delta\psi^\ast.
\end{equation}
The moving wall enters the linearized GNLS problem through \eqref{eq:part01-delta-V-conf}.  In BdG form the schematic system is
\begin{equation}
\label{eq:part01-BdG-schematic}
i\hbar\partial_t
\begin{pmatrix}
\delta\psi\\ \delta\psi^\ast
\end{pmatrix}
=
\mathcal L_{\rm BdG}
\begin{pmatrix}
\delta\psi\\ \delta\psi^\ast
\end{pmatrix}
+\mathcal C_A[\delta A_M]+\mathcal C_\eta[\eta].
\end{equation}
Stage~003 does not claim to solve the full BdG PDE.  It passes to a stable normal-mode reduction in selected \(l=0\) and \(l=2\) sectors.  Once that reduction is declared, integrating out the stable modes is exact algebra.

\subsection{Axisymmetric BdG support kernel}
\label{subsec:part01-axisymmetric-BdG-kernel}

Let \(Q^A=(\delta a,\delta L)\) and let \(X_\alpha\) be stable scalar BdG normal coordinates with frequencies \(\varpi_{0\alpha}>0\).  The reduced quadratic Lagrangian is
\begin{equation}
\label{eq:part01-Lred0-BdG}
L_{\rm red}^{(0)}
=
\frac12\dot Q^TM_0\dot Q-\frac12Q^TK_0Q
+
\frac12\sum_\alpha\left(\dot X_\alpha^2-\varpi_{0\alpha}^2X_\alpha^2\right)
+
\sum_{A,\alpha}C_{A\alpha}Q^AX_\alpha.
\end{equation}
In frequency space, the matter coordinates solve
\begin{equation}
(\varpi_{0\alpha}^2-\omega^2)X_\alpha=C_{A\alpha}Q^A.
\end{equation}
Substitution gives the exact effective matrix kernel
\begin{equation}
\label{eq:part01-D0eff-BdG}
\boxed{
D_0^{\rm eff}(\omega)
=
K_0-\omega^2M_0-C(\Omega_0^2-\omega^2I)^{-1}C^T},
\end{equation}
where
\begin{equation}
\Omega_0^2=\operatorname{diag}(\varpi_{0\alpha}^2).
\end{equation}
Its low-frequency expansion is
\begin{equation}
D_0^{\rm eff}(\omega)
=
K_0^{\rm eff}-\omega^2M_0^{\rm eff}-\omega^4N_0^{\rm eff}+O(\omega^6),
\end{equation}
with
\begin{equation}
\label{eq:part01-BdG-axisymmetric-moments}
\boxed{
K_0^{\rm eff}=K_0-C\Omega_0^{-2}C^T,
\qquad
M_0^{\rm eff}=M_0+C\Omega_0^{-4}C^T,
\qquad
N_0^{\rm eff}=C\Omega_0^{-6}C^T}.
\end{equation}
Thus stable matter support softens the static stiffness, increases the effective inertia, and generates the next even low-frequency coefficient.

\subsection{Grouped real \texorpdfstring{\(P_2\)}{P2} BdG kernels}
\label{subsec:part01-p2-BdG-kernels}

For each grouped lane \(A\in\{20,21,22\}\), let \(q_A\) be the wall amplitude and \(X_{A\alpha}\) be stable quadrupole BdG coordinates with frequencies \(\varpi_{A\alpha}\).  The reduced Lagrangian is
\begin{equation}
L_{\rm red}^{(2)}
=
\sum_A\left[
\frac12M_A\dot q_A^2-\frac12K_Aq_A^2
+
\frac12\sum_\alpha(\dot X_{A\alpha}^2-\varpi_{A\alpha}^2X_{A\alpha}^2)
+
\sum_\alpha g_{A\alpha}q_AX_{A\alpha}
\right].
\end{equation}
Eliminating the stable coordinates gives
\begin{equation}
\label{eq:part01-DAeff-BdG}
\boxed{
D_A^{\rm eff}(\omega)
=
K_A-M_A\omega^2
-
\sum_\alpha\frac{g_{A\alpha}^2}{\varpi_{A\alpha}^2-\omega^2}}.
\end{equation}
The low-frequency coefficients are
\begin{equation}
\label{eq:part01-BdG-P2-coefficients}
\boxed{
\begin{aligned}
d_{0A}^{\rm eff}&=K_A-\sum_\alpha\frac{g_{A\alpha}^2}{\varpi_{A\alpha}^2},\\
d_{2A}^{\rm eff}&=-\left(M_A+\sum_\alpha\frac{g_{A\alpha}^2}{\varpi_{A\alpha}^4}\right),\\
d_{4A}^{\rm eff}&=-\sum_\alpha\frac{g_{A\alpha}^2}{\varpi_{A\alpha}^6}.
\end{aligned}
}
\end{equation}
These are the first explicit conservative grouped-lane moments generated by matter support.

\subsection{First pole-shift formula}
\label{subsec:part01-first-pole-shift}

For one wall mode \(q\) coupled to one stable matter mode \(X\),
\begin{equation}
L
=
\frac12M\dot q^2-\frac12Kq^2+\frac12\dot X^2-\frac12\varpi^2X^2+gqX.
\end{equation}
The exact dispersion relation is
\begin{equation}
\label{eq:part01-pole-dispersion}
(K-M\omega^2)(\varpi^2-\omega^2)-g^2=0.
\end{equation}
With \(\Omega_\eta^2=K/M\), the exact poles are
\begin{equation}
\label{eq:part01-pole-shifts-exact}
\boxed{
\omega_\pm^2
=
\frac{\Omega_\eta^2+\varpi^2
\pm\sqrt{(\Omega_\eta^2-\varpi^2)^2+4g^2/M}}{2}}.
\end{equation}
If \(\varpi^2>\Omega_\eta^2\) and \(g\) is weak, the wall-like pole shifts by
\begin{equation}
\label{eq:part01-wall-pole-shift}
\boxed{
\delta\Omega_\eta^2
=
-\frac{g^2/M}{\varpi^2-\Omega_\eta^2}+O(g^4)}.
\end{equation}
This is the first explicit reduced origin of the pole data that later appear as open throat observables.

\subsection{Matter-coupled isotropy diagnostic}
\label{subsec:part01-BdG-isotropy-diagnostic}

For grouped triples \((x_{20},x_{21},x_{22})\), define
\begin{equation}
\label{eq:part01-grouped-bar-a-b-early}
\bar x=\frac{x_{20}+2x_{21}+2x_{22}}5,
\qquad
a_x=\frac{2x_{20}-x_{21}-x_{22}}{10},
\qquad
b_x=\frac{x_{21}-x_{22}}2.
\end{equation}
If
\begin{equation}
K_{20}=K_{21}=K_{22},
\quad
M_{20}=M_{21}=M_{22},
\quad
g_{20,\alpha}=g_{21,\alpha}=g_{22,\alpha},
\quad
\varpi_{20,\alpha}=\varpi_{21,\alpha}=\varpi_{22,\alpha},
\end{equation}
then all grouped coefficients in \eqref{eq:part01-BdG-P2-coefficients} are lane-independent, and therefore the grouped anisotropy defects vanish.

\begin{theorem}[BdG Schur-complement support kernel]
\label{thm:part01-BdG-schur}
Inside the stable normal-mode reduction of the linearized matter sector, integrating out the BdG coordinates gives the exact Schur-complement kernels \eqref{eq:part01-D0eff-BdG} and \eqref{eq:part01-DAeff-BdG}.  The low-frequency moments are fixed by the support frequencies and coupling overlaps; grouped \(P_2\) anisotropy can enter only through lane dependence of the wall data, support frequencies, or coupling overlaps.
\end{theorem}

\begin{proof}
The reduced matter coordinates enter quadratically and linearly.  Their frequency-domain equations are algebraic.  Solving those equations and substituting back into the wall equations gives the Schur complements.  Expanding \((\Omega^2-\omega^2)^{-1}\) in powers of \(\omega^2\) gives the stated low-frequency coefficients.  The isotropy statement follows from equality of the three lane data.
\end{proof}

\section{Projection-first Maxwell/mixed response and outgoing \texorpdfstring{\(l=2\)}{l=2} bridge}
\label{sec:part01_parent_geometry:localized-maxwell-mixed-response-and-outgoing-l-2-}

\subsection{Why the mixed sector enters before the outgoing bridge}
\label{subsec:part01-mixed-bridge-purpose}

The BdG support kernel is conservative.  It can produce even low-frequency response moments and pole shifts, but by itself it does not produce an honest passive/outgoing odd term.  The first available microscopic carrier of an outgoing branch is the localized Maxwell/mixed block, because \(A_w\), \(F_{\mu w}\), and \(J^w\) are present in the parent ontology and are not gauge artifacts.

Stages~004--021 first project the localized parent Maxwell equation, producing the exact open-system law \eqref{eq:part01-projected-maxwell-law}.  Only after the matched channel \eqref{eq:part01-projection-reduction-match} is declared does the derivation compress the EM sector into a minimal reduced lane containing one wall coordinate \(Q_l\), one brane-like localized Maxwell coordinate \(A_l\), and one mixed-sector coordinate \(W_l\).  This is not the full PDE.  It is the Stage~021 normal form used to carry both conservative gauge self-energy and passive outgoing transfer into the grouped bundle.

\subsection{One-lane Maxwell/mixed Lagrangian}
\label{subsec:part01-maxwell-mixed-L}

The reduced quadratic Lagrangian is
\begin{equation}
\label{eq:part01-maxwell-mixed-L}
\begin{aligned}
L_l={}&\frac12M_l\dot Q_l^2-\frac12K_lQ_l^2
+\frac12\dot A_l^2-\frac12\Omega_{A,l}^2A_l^2
+\frac12\dot W_l^2-\frac12\Omega_{W,l}^2W_l^2 \\
&+R_lA_lW_l+g_{A,l}Q_lA_l+g_{W,l}Q_lW_l .
\end{aligned}
\end{equation}
Here \(A_l\) is a brane-like localized gauge coordinate, \(W_l\) is a mixed \(A_w/F_{\mu w}/J^w\)-active coordinate, \(R_l\) mixes the internal gauge coordinates, and \(g_{A,l},g_{W,l}\) are wall couplings.

\subsection{Conservative self-energy}
\label{subsec:part01-conservative-maxwell-self-energy}

With frequency convention \(e^{-i\omega t}\), define
\begin{equation}
A_l(\omega)=\Omega_{A,l}^2-\omega^2,
\qquad
W_l(\omega)=\Omega_{W,l}^2-\omega^2,
\qquad
\Delta_l(\omega)=A_l(\omega)W_l(\omega)-R_l^2.
\end{equation}
Eliminating \(A_l,W_l\) gives the conservative self-energy
\begin{equation}
\label{eq:part01-sigma-cons}
\boxed{
\Sigma_l^{\rm EM+mix,cons}(\omega)
=
\frac{g_{A,l}^2W_l(\omega)+2g_{A,l}g_{W,l}R_l+g_{W,l}^2A_l(\omega)}
{\Delta_l(\omega)}}.
\end{equation}
The wall operator is
\begin{equation}
D_l^{\rm cons}(\omega)=K_l-M_l\omega^2-\Sigma_l^{\rm EM+mix,cons}(\omega).
\end{equation}
Set
\begin{equation}
\Delta_{0,l}=\Omega_{A,l}^2\Omega_{W,l}^2-R_l^2,
\qquad
S_{2,l}=\Omega_{A,l}^2+\Omega_{W,l}^2,
\end{equation}
\begin{equation}
N_{0,l}=g_{A,l}^2\Omega_{W,l}^2+2g_{A,l}g_{W,l}R_l+g_{W,l}^2\Omega_{A,l}^2,
\qquad
G_{2,l}=g_{A,l}^2+g_{W,l}^2.
\end{equation}
Then
\begin{equation}
\Sigma_l^{\rm EM+mix,cons}(\omega)
=
z_{l,0}+z_{l,2}\omega^2+z_{l,4}\omega^4+O(\omega^6),
\end{equation}
with
\begin{equation}
\label{eq:part01-z-coeffs}
\boxed{
\begin{aligned}
z_{l,0}&=\frac{N_{0,l}}{\Delta_{0,l}},\\
z_{l,2}&=\frac{N_{0,l}S_{2,l}-G_{2,l}\Delta_{0,l}}{\Delta_{0,l}^2},\\
z_{l,4}&=\frac{N_{0,l}(S_{2,l}^2-\Delta_{0,l})-S_{2,l}G_{2,l}\Delta_{0,l}}{\Delta_{0,l}^3}.
\end{aligned}
}
\end{equation}

\subsection{Outgoing-port dressing and transfer factor}
\label{subsec:part01-outgoing-port-dressing}

Attach the exterior passive/outgoing channel to the mixed coordinate by replacing
\begin{equation}
W_l(\omega)\longrightarrow W_l(\omega)-\Pi_l^{\rm out}(\omega).
\end{equation}
To first order in the outgoing port,
\begin{equation}
\Sigma_l^{\rm full}(\omega)
=
\Sigma_l^{\rm cons}(\omega)+\Pi_l^{\rm out}(\omega)N_l(\omega)
+
O\!\left((\Pi_l^{\rm out})^2\right),
\end{equation}
where the transfer factor is
\begin{equation}
\label{eq:part01-outgoing-transfer-N}
\boxed{
N_l(\omega)
=
\frac{\bigl[A_l(\omega)g_{W,l}+R_lg_{A,l}\bigr]^2}
{\bigl[A_l(\omega)W_l(\omega)-R_l^2\bigr]^2}}.
\end{equation}
At zero frequency,
\begin{equation}
\label{eq:part01-N-zero-positive}
\boxed{
N_l(0)
=
\frac{\bigl[\Omega_{A,l}^2g_{W,l}+R_lg_{A,l}\bigr]^2}
{\bigl[\Omega_{A,l}^2\Omega_{W,l}^2-R_l^2\bigr]^2}
\ge0}.
\end{equation}
The sign of the outgoing odd term is therefore supplied by the outgoing port; the internal conservative block transfers it with a nonnegative static factor.

\subsection{Compact outgoing \texorpdfstring{\(l=2\)}{l=2} fingerprint}
\label{subsec:part01-l2-outgoing-fingerprint}

For the compact outgoing \(l=2\) branch, the normalized outgoing response has the low-frequency form
\begin{equation}
\label{eq:part01-l2-fingerprint}
\widehat Y_2^{\rm out}(\omega)
=
1+\frac{a^2\omega^2}{9c_s^2}
+\frac{4a^4\omega^4}{81c_s^4}
+i\frac{a^5\omega^5}{27c_s^5}
+O(\omega^6).
\end{equation}
Thus the outgoing port coefficient is
\begin{equation}
\label{eq:part01-Gamma2-port}
\Gamma_2^{\rm port}=\frac{a^5}{27c_s^5}.
\end{equation}
Combining \eqref{eq:part01-outgoing-transfer-N} with \eqref{eq:part01-Gamma2-port}, the wall-level odd quadrupole contribution is
\begin{equation}
\label{eq:part01-wall-odd-quadrupole}
\boxed{
\delta D_2^{\rm odd}(\omega)
=
-iN_2(0)\frac{a^5}{27c_s^5}\omega^5+O(\omega^7)}.
\end{equation}
In the wall-operator convention \(D=K-M\omega^2-\Sigma\), this is the passive sign.  In normalized response language the same branch is read with the positive \(+i\omega^5\) fingerprint.

\subsection{Scalar compatibility check}
\label{subsec:part01-scalar-compatibility}

A naive scalar outlet could introduce a dangerous \(i\omega\) monopole term.  The reduced mixed-sector model shows how such a term can be delayed.  If the scalar/breathing wall coordinate couples to the port-active mixed combination only through a derivative coupling, for example
\begin{equation}
g_{W,0}(\omega)=\eta\omega,
\end{equation}
and if there is no direct non-derivative coupling into the same port-active combination, then
\begin{equation}
N_0(\omega)=O(\omega^2).
\end{equation}
Therefore an outgoing scalar port with \(\Pi_0^{\rm out}(\omega)=i\gamma_1\omega+\cdots\) contributes at wall level only as \(O(i\omega^3)\).  This is a compatibility mechanism, not yet a full scalar no-go theorem.

\begin{theorem}[Projection-first Maxwell/mixed outgoing bridge]
\label{thm:part01-maxwell-mixed-bridge}
The exact projected Maxwell law \eqref{eq:part01-projected-maxwell-law} retains the transverse mixed channel before any brane reduction.  In the matched one-lane Maxwell/mixed normal form, the wall inherits an outgoing odd contribution through the exact transfer factor \eqref{eq:part01-outgoing-transfer-N}.  For \(l=2\), the leading wall-level odd coefficient is \eqref{eq:part01-wall-odd-quadrupole}; the remaining normalization problem is the computation of the static transfer factor on the actual moving-throat branch.
\end{theorem}

\begin{proof}
Projection supplies the mixed transverse source term in \eqref{eq:part01-projected-maxwell-law}; the matched channel \eqref{eq:part01-projection-reduction-match} is the reduction in which it can be represented by a finite one-port packet.  Eliminating the two internal gauge coordinates from \eqref{eq:part01-maxwell-mixed-L} gives the conservative self-energy \eqref{eq:part01-sigma-cons}.  Replacing \(W_l\) by \(W_l-\Pi_l^{\rm out}\) and expanding to first order in \(\Pi_l^{\rm out}\) gives \eqref{eq:part01-outgoing-transfer-N}.  Multiplying by the compact outgoing \(l=2\) fingerprint coefficient \eqref{eq:part01-Gamma2-port} gives \eqref{eq:part01-wall-odd-quadrupole}.
\end{proof}

\section{Full grouped real \texorpdfstring{\(P_2\)}{P2} bundle transition}
\label{sec:part01_parent_geometry:full-grouped-real-p-2-bundle-transition}

\subsection{Conservative operator moments and normalized response moments}
\label{subsec:part01-operator-to-response}

Stages~022 and 023 lift the projection-first Maxwell packet, through its matched one-port normal form, into the full grouped real bundle.  For each grouped lane \(A\in\{20,21,22\}\), write the conservative operator as
\begin{equation}
\label{eq:part01-DA-cons-expansion}
D_A^{\rm cons}(\omega)
=
D_{A0}+D_{A2}\omega^2+D_{A4}\omega^4+O(\omega^6).
\end{equation}
The normalized response is
\begin{equation}
\label{eq:part01-Y-A-response}
Y_A(\omega)
=
\frac{D_{A0}}{D_A^{\rm cons}(\omega)}
=
1+u_2^{(A)}\omega^2+u_4^{(A)}\omega^4+O(\omega^6).
\end{equation}
The exact conversion formulas are
\begin{equation}
\label{eq:part01-u2-u4-conversion}
\boxed{
u_2^{(A)}=-\frac{D_{A2}}{D_{A0}},
\qquad
u_4^{(A)}=\frac{D_{A2}^2-D_{A0}D_{A4}}{D_{A0}^2}}.
\end{equation}
This is the first bookkeeping bridge between the wall/BdG/Maxwell operator language and the grouped-response language used by the later retarded quadrupole tests.

\subsection{Grouped inverse map and weighted projectors}
\label{subsec:part01-weighted-projectors}

For a grouped triple \(x=(x_{20},x_{21},x_{22})^T\), define
\begin{equation}
\label{eq:part01-bar-a-b}
\bar x=\frac{x_{20}+2x_{21}+2x_{22}}5,
\qquad
a_x=\frac{2x_{20}-x_{21}-x_{22}}{10},
\qquad
b_x=\frac{x_{21}-x_{22}}2.
\end{equation}
The inverse map is
\begin{equation}
\label{eq:part01-inverse-bar-a-b}
\boxed{
x_{20}=\bar x+4a_x,
\qquad
x_{21}=\bar x-a_x+b_x,
\qquad
x_{22}=\bar x-a_x-b_x}.
\end{equation}
The natural grouped metric is
\begin{equation}
G_{\rm grp}=\operatorname{diag}(1,2,2).
\end{equation}
A \(G_{\rm grp}\)-orthogonal basis is
\begin{equation}
e_{\bar x}=(1,1,1),
\qquad
e_a=(4,-1,-1),
\qquad
e_b=(0,1,-1),
\end{equation}
with norms
\begin{equation}
e_{\bar x}^TG_{\rm grp}e_{\bar x}=5,
\qquad
e_a^TG_{\rm grp}e_a=20,
\qquad
e_b^TG_{\rm grp}e_b=4.
\end{equation}
The exact projectors are
\begin{equation}
\label{eq:part01-grouped-projectors}
\boxed{
P_{\bar x}=\frac15
\begin{pmatrix}
1&2&2\\1&2&2\\1&2&2
\end{pmatrix},
\quad
P_a=\frac1{20}
\begin{pmatrix}
16&-8&-8\\-4&2&2\\-4&2&2
\end{pmatrix},
\quad
P_b=\frac14
\begin{pmatrix}
0&0&0\\0&2&-2\\0&-2&2
\end{pmatrix}}.
\end{equation}
They satisfy
\begin{equation}
P_{\bar x}+P_a+P_b=I_3,
\qquad
P_iP_j=\delta_{ij}P_i.
\end{equation}
The grouped isotropy condition is simply
\begin{equation}
\label{eq:part01-isotropy-projector-condition}
P_ax=P_bx=0.
\end{equation}

\subsection{Outgoing prefactor and normalized branch coefficients}
\label{subsec:part01-outgoing-prefactor}

Let the outgoing transfer factor in lane \(A\) be
\begin{equation}
N_A(\omega)=N_{A0}+N_{A2}\omega^2+N_{A4}\omega^4+O(\omega^6).
\end{equation}
The internal prefactor multiplying the outgoing branch is
\begin{equation}
\label{eq:part01-PrefA-def}
\operatorname{Pref}_A(\omega)
=
\frac{D_{A0}N_A(\omega)}{D_A^{\rm cons}(\omega)^2}
=
P_{A0}+P_{A2}\omega^2+P_{A4}\omega^4+O(\omega^6).
\end{equation}
The exact coefficients are
\begin{equation}
\label{eq:part01-prefactor-coeffs}
\boxed{
\begin{aligned}
P_{A0}&=\frac{N_{A0}}{D_{A0}},\\
P_{A2}&=\frac{D_{A0}N_{A2}-2D_{A2}N_{A0}}{D_{A0}^2},\\
P_{A4}&=
\frac{D_{A0}^2N_{A4}-2D_{A0}(D_{A2}N_{A2}+D_{A4}N_{A0})+3D_{A2}^2N_{A0}}
{D_{A0}^3}.
\end{aligned}
}
\end{equation}
Multiplying by the compact outgoing \(l=2\) fingerprint \eqref{eq:part01-l2-fingerprint}, with
\begin{equation}
A_2^{\rm out}=\frac{a^2}{9c_s^2},
\qquad
B_2^{\rm out}=\frac{4a^4}{81c_s^4},
\qquad
G_5^{\rm out}=\frac{a^5}{27c_s^5},
\end{equation}
gives
\begin{equation}
\label{eq:part01-branch-K-Gamma}
\boxed{
K_{A0}=P_{A0},
\quad
K_{A2}=P_{A2}+A_2^{\rm out}P_{A0},
\quad
K_{A4}=P_{A4}+A_2^{\rm out}P_{A2}+B_2^{\rm out}P_{A0},
\quad
\Gamma_5^{(A)}=G_5^{\rm out}P_{A0}}.
\end{equation}
Only the static prefactor \(P_{A0}\) enters the leading \(i\omega^5\) coefficient.

\subsection{Full reduced grouped bundle}
\label{subsec:part01-full-bundle}

Stage~023 writes the full coupled lane data explicitly.  For each lane \(A\), define BdG support moments
\begin{equation}
\label{eq:part01-B-moments}
B_{A0}=\sum_\alpha\frac{c_{A\alpha}^2}{\varpi_{A\alpha}^2},
\qquad
B_{A2}=\sum_\alpha\frac{c_{A\alpha}^2}{\varpi_{A\alpha}^4},
\qquad
B_{A4}=\sum_\alpha\frac{c_{A\alpha}^2}{\varpi_{A\alpha}^6}.
\end{equation}
For each conservative Maxwell/mixed internal pair \(r\), set
\begin{equation}
\Delta_{A,r}=\Omega_{U,A,r}^2\Omega_{W,A,r}^2-R_{A,r}^2,
\qquad
S_{A,r}=\Omega_{U,A,r}^2+\Omega_{W,A,r}^2,
\end{equation}
\begin{equation}
Q_{A,r}
=
g_{U,A,r}^2\Omega_{W,A,r}^2
+
2g_{U,A,r}g_{W,A,r}R_{A,r}
+
g_{W,A,r}^2\Omega_{U,A,r}^2,
\qquad
G_{A,r}=g_{U,A,r}^2+g_{W,A,r}^2.
\end{equation}
Then
\begin{equation}
\label{eq:part01-Z-moments}
\boxed{
\begin{aligned}
Z_{A0}^{(r)}&=\frac{Q_{A,r}}{\Delta_{A,r}},\\
Z_{A2}^{(r)}&=\frac{Q_{A,r}S_{A,r}-G_{A,r}\Delta_{A,r}}{\Delta_{A,r}^2},\\
Z_{A4}^{(r)}&=
\frac{Q_{A,r}(S_{A,r}^2-\Delta_{A,r})-S_{A,r}G_{A,r}\Delta_{A,r}}
{\Delta_{A,r}^3}.
\end{aligned}
}
\end{equation}
The outgoing-transfer moments are
\begin{equation}
\label{eq:part01-N-moments}
\boxed{
\begin{aligned}
N_{A0}^{(r)}&=\frac{P_{A,r}^2}{\Delta_{A,r}^2},\\
N_{A2}^{(r)}&=
\frac{2P_{A,r}(P_{A,r}S_{A,r}-\Delta_{A,r}g_{W,A,r})}
{\Delta_{A,r}^3},\\
N_{A4}^{(r)}&=
\frac{\Delta_{A,r}^2g_{W,A,r}^2-2\Delta_{A,r}P_{A,r}^2-4\Delta_{A,r}P_{A,r}S_{A,r}g_{W,A,r}+3P_{A,r}^2S_{A,r}^2}
{\Delta_{A,r}^4},
\end{aligned}
}
\end{equation}
where
\begin{equation}
P_{A,r}=\Omega_{U,A,r}^2g_{W,A,r}+R_{A,r}g_{U,A,r}.
\end{equation}
Summing over \(r\) gives \(Z_{An}\) and \(N_{An}\).  The total conservative wall coefficients are
\begin{equation}
\label{eq:part01-full-D-coefficients}
\boxed{
D_{A0}=K_A-B_{A0}-Z_{A0},
\qquad
D_{A2}=-\bigl(M_A+B_{A2}+Z_{A2}\bigr),
\qquad
D_{A4}=-\bigl(B_{A4}+Z_{A4}\bigr)}.
\end{equation}

\subsection{Projected Maxwell slippage into grouped slots}
\label{subsec:part01-projected-maxwell-slippage}

The projection-first correction is transported into the same grouped variables
by the slot shifts
\begin{equation}
\label{eq:part01-projected-slot-shifts}
Z_n\mapsto Z_n+\varepsilon z_n,
\qquad
N_n\mapsto N_n+\varepsilon n_n,
\qquad
n\in\{0,2,4\}.
\end{equation}
At a symmetric interior slice \(\varepsilon=O(\sigma^2)\), while a one-sided
mouth projection gives \(\varepsilon=O(\ell)\).  On an isotropic lane, the
induced response shifts are
\begin{equation}
\label{eq:part01-projected-du2-du4}
\delta u_2=\frac{D_0z_2-D_2z_0}{D_0^2},
\qquad
\delta u_4=
\frac{D_0^2z_4-D_0(2D_2z_2+D_4z_0)+2D_2^2z_0}{D_0^3},
\end{equation}
and the static outgoing-prefactor shift is
\begin{equation}
\label{eq:part01-projected-dp0}
\delta P_0=\frac{D_0n_0+N_0z_0}{D_0^2}.
\end{equation}
For the eliminated isotropic one-pole compatibility surface
\begin{equation}
\mathcal C=\frac{N_0}{P_{0,\rm target}}-\frac{3S^2}{T},
\qquad
S=M+B_2+Z_2,\quad T=B_4+Z_4,
\end{equation}
the projected variation is
\begin{equation}
\label{eq:part01-projected-compatibility-variation}
\boxed{
\delta\mathcal C
=\frac{n_0}{P_{0,\rm target}}
-\frac{6Sz_2}{T}
+\frac{3S^2z_4}{T^2}}.
\end{equation}
The static conservative shift \(z_0\) cancels from this eliminated surface.
It still changes \(P_0\) through \eqref{eq:part01-projected-dp0}, so the
projection-first EM correction remains live in the normalization branch even
when it does not move the one-pole compatibility equation.

\subsection{Isotropic branch and final Part I normalization test}
\label{subsec:part01-isotropic-normalization-test}

On the isotropic branch,
\begin{equation}
D_{20,n}=D_{21,n}=D_{22,n}=D_n,
\qquad
N_{20,n}=N_{21,n}=N_{22,n}=N_n.
\end{equation}
Then
\begin{equation}
\label{eq:part01-isotropic-u-P}
\boxed{
\begin{aligned}
u_2&=-\frac{D_2}{D_0},\\
u_4&=\frac{D_2^2-D_0D_4}{D_0^2},\\
P_0&=\frac{N_0}{D_0},\\
P_2&=\frac{D_0N_2-2D_2N_0}{D_0^2},\\
P_4&=
\frac{D_0^2N_4-2D_0(D_2N_2+D_4N_0)+3D_2^2N_0}
{D_0^3}.
\end{aligned}
}
\end{equation}
The universal quadrupole-normalization condition is
\begin{equation}
\label{eq:part01-normalization-condition}
\boxed{
\widehat m_0^{\,2}P_0
=
\widehat m_0^{\,2}\frac{N_0}{D_0}
=
\frac{54G c_s^5}{5a^5c^5}}.
\end{equation}
Using \eqref{eq:part01-full-D-coefficients}, the same test is
\begin{equation}
\label{eq:part01-final-normalization}
\boxed{
\widehat m_0^{\,2}
\frac{N_0}{K-B_0-Z_0}
=
\frac{54G c_s^5}{5a^5c^5}}.
\end{equation}
This is the central output of Part I.  The algebraic machinery has compressed the outgoing-normalization issue to a single isotropic ratio.  The actual branch still has to compute the numerator, denominator, and source-map factor.

\subsection{Constant-prefactor branch}
\label{subsec:part01-constant-prefactor-branch}

The constant-prefactor branch used in the later retarded program is defined by
\begin{equation}
P_2=0,
\qquad
P_4=0.
\end{equation}
Using \eqref{eq:part01-isotropic-u-P}, this is equivalent to the exact correlation conditions
\begin{equation}
\label{eq:part01-constant-prefactor-conditions}
\boxed{
N_2=\frac{2D_2N_0}{D_0},
\qquad
N_4=\frac{2D_0(D_2N_2+D_4N_0)-3D_2^2N_0}{D_0^2}}.
\end{equation}
Thus the constant-prefactor branch does not require \(N_2\) and \(N_4\) to vanish individually.  It requires them to be locked to the conservative moments \(D_0,D_2,D_4\).

\subsection{First-order anisotropy transport}
\label{subsec:part01-anisotropy-transport}

Linearize around an isotropic branch:
\begin{equation}
D_{A,n}=D_n+\delta D_{A,n},
\qquad
N_{A,n}=N_n+\delta N_{A,n}.
\end{equation}
The response-moment variation is
\begin{equation}
\delta u_2^{(A)}
=
-\frac{\delta D_{A,2}+u_2\delta D_{A,0}}{D_0}.
\end{equation}
Therefore
\begin{equation}
\label{eq:part01-u-anisotropy-transport}
\boxed{
a_{u,2}=-\frac{a_{D,2}+u_2a_{D,0}}{D_0},
\qquad
b_{u,2}=-\frac{b_{D,2}+u_2b_{D,0}}{D_0}}.
\end{equation}
Similarly, since \(P_0^{(A)}=N_{A0}/D_{A0}\),
\begin{equation}
\delta P_0^{(A)}
=
\frac{\delta N_{A,0}-P_0\delta D_{A,0}}{D_0},
\end{equation}
and hence
\begin{equation}
\label{eq:part01-P0-anisotropy-transport}
\boxed{
a_{P,0}=\frac{a_{N,0}-P_0a_{D,0}}{D_0},
\qquad
b_{P,0}=\frac{b_{N,0}-P_0b_{D,0}}{D_0}}.
\end{equation}
These formulas will be used later when weak-axisymmetric anisotropy is transported through the grouped \(P_2\) bundle.

\begin{theorem}[Full grouped \(P_2\) bundle and isotropic normalization test]
\label{thm:part01-full-grouped-bundle}
Inside the full reduced wall/BdG/projected-Maxwell grouped-bundle closure, the coupled conservative moments and outgoing-transfer moments are given by \eqref{eq:part01-B-moments}--\eqref{eq:part01-full-D-coefficients}, with projection-first EM slippage transported by \eqref{eq:part01-projected-slot-shifts}--\eqref{eq:part01-projected-compatibility-variation}.  On the isotropic branch the leading outgoing quadrupole normalization reduces to the single scalar condition \eqref{eq:part01-final-normalization}.  First-order anisotropy around that branch is transported by \eqref{eq:part01-u-anisotropy-transport} and \eqref{eq:part01-P0-anisotropy-transport}.
\end{theorem}

\begin{proof}
The full grouped bundle is the direct sum of the lane-wise reduced mechanisms from Stage~003 and Stages~004--021.  The BdG support moments, conservative Maxwell/mixed moments, and outgoing-transfer moments add to the wall operator as shown in \eqref{eq:part01-full-D-coefficients}; projection-first corrections enter the same \(Z_n,N_n\) slots by \eqref{eq:part01-projected-slot-shifts}.  The normalized response and prefactor expansions are algebraic inversions of \(D_A^{\rm cons}\) and \(D_A^{\rm cons}{}^2\), giving \eqref{eq:part01-isotropic-u-P}.  Equality of the three grouped lanes kills the projector components \(P_a\) and \(P_b\), leaving the scalar ratio \(N_0/D_0\).  Linearizing the definitions of \(u_2^{(A)}\) and \(P_0^{(A)}\) gives the anisotropy transport laws.
\end{proof}

\section{Part I audit summary and handoff to Part II}
\label{sec:part01-audit-summary}

Part I establishes the following stable outputs.
\begin{enumerate}[leftmargin=2.2em]
  \item \textbf{Exact parent imports.}  The gauged GNLS equation, localized Maxwell equation, current conservation, projection language, charge firewall, and mixed-sector gauge invariants are fixed by the declared parent theory.
  \item \textbf{Moving geometry closure.}  The working moving-throat field is \(\Sigma=r-R(\Omega,w,t)\), with \(a(t)\) and \(L(t)\) recovered as collective moments.
  \item \textbf{Wall modal structure.}  The distributed wall action separates the scalar lane, residual geometry lane, grouped real \(P_2\) lane, and higher multipoles.
  \item \textbf{Breathing reduction.}  The old \((a,L)\) closure is the lowest axisymmetric truncation of the distributed wall field.
  \item \textbf{Matter support Schur complement.}  Stable BdG modes generate conservative low-frequency support moments and pole shifts.
  \item \textbf{Projected Maxwell/mixed outgoing bridge.}  Projection-first Maxwell supplies the open-system mixed-channel provenance; the matched one-port normal form supplies a nonnegative outgoing-transfer factor and carries the compact outgoing \(l=2\) \(i\omega^5\) fingerprint.
  \item \textbf{Grouped bundle.}  The final Part I object is the grouped coefficient packet \((D_A,N_A,P_A)\) and the isotropic normalization gate \eqref{eq:part01-final-normalization}.
\end{enumerate}

This chapter does not prove that the actual nonlinear moving-throat PDE realizes the target value in \eqref{eq:part01-final-normalization}.  It also does not prove uniqueness of the shape-field closure, solve the full BdG spectrum, solve the full Maxwell/mixed spectrum, or complete the passive/outgoing boundary-value problem.  Those remain branch-realization tasks.

\begin{verificationchecklist}
\item The parent GNLS, Maxwell, continuity, and mixed-sector equations are imported with corrected charge ontology.
\item The sign in \(\delta V_{\rm conf}=-V'_{\rm wall}\eta/\ell_c\) is preserved.
\item The monopole normalization \(q_{00}=2\sqrt\pi\,\delta a\) is preserved.
\item The old \((a,L)\) closure is recovered as a two-mode axisymmetric truncation.
\item The geometry-only grouped real \(P_2\) sector is degenerate on an isotropic reference throat.
\item Stable BdG modes are integrated out by Schur complement, not by a fitted response coefficient.
\item The projection-first localized Maxwell block retains \(A_w,F_{\mu w},J^w\) until a controlled matched brane reduction is declared.
\item The outgoing \(l=2\) odd coefficient depends on the static prefactor \(P_0=N_0/D_0\).
\item The grouped metric \(G_{\rm grp}=\operatorname{diag}(1,2,2)\) and projectors \(P_{\bar x},P_a,P_b\) are used consistently.
\item The final Part I theorem gate is the branch-realization test \eqref{eq:part01-final-normalization}, not a free normalization choice.
\end{verificationchecklist}

Part II uses the operator-to-response bridge \eqref{eq:part01-u2-u4-conversion} and the grouped projector calculus to build the conservative response and overlap-isotropy tests.  Part III uses the same full-bundle data to discuss support completion and continuum placement.  Part IV uses the outgoing bridge and scalar compatibility mechanism when the retarded closure, outlet models, core/mouth compensation, and boundary-layer laws are introduced.  Part V uses the anisotropy transport formulas \eqref{eq:part01-u-anisotropy-transport} and \eqref{eq:part01-P0-anisotropy-transport} as the starting point for the weak-axisymmetric quotient and orbit-lock structure.

The citation-safe summary is:
\begin{quote}
Stages~001--023 define the moving-throat linearized response backbone.  They project the parent Maxwell sector before applying the matched one-port reduction, and reduce the parent wall/matter/gauge system to a grouped real \(P_2\) response bundle with exact conservative moments, exact outgoing-transfer moments, exact grouped projectors, and the isotropic normalization test
\[
\widehat m_0^{\,2}\frac{N_0}{K-B_0-Z_0}=\frac{54Gc_s^5}{5a^5c^5}.
\]
They do not by themselves prove that the actual nonlinear branch realizes that test.
\end{quote}
